\documentclass[aspectratio=169,11pt]{beamer}

\usetheme[
  progressbar=frametitle,
  numbering=fraction
]{metropolis}

\usepackage{ctex}
\usepackage{booktabs}
\usepackage{graphicx}

\title{用 Beamer 做一份清爽的学术汇报}
\subtitle{从主题选择到页面结构}
\author{戚柒}
\institute{Personal Blog Demo}
\date{\today}

\begin{document}

\maketitle

\begin{frame}{今天要解决的问题}
  \begin{itemize}
    \item 找到适合学术和技术汇报的 Beamer 主题
    \item 用 Metropolis 快速搭建一份 16:9 幻灯片
    \item 把内容拆成清晰的章节、观点和图表
  \end{itemize}
\end{frame}

\section{为什么选 Metropolis}

\begin{frame}{主题选择标准}
  \begin{columns}[T,onlytextwidth]
    \column{0.48\textwidth}
      \textbf{适合长期使用}
      \begin{itemize}
        \item 视觉噪音少
        \item 默认排版稳
        \item 适合公式和代码
      \end{itemize}

    \column{0.48\textwidth}
      \textbf{容易交给别人}
      \begin{itemize}
        \item CTAN/TeX Live 常见
        \item Overleaf 可直接试
        \item 只需少量配置
      \end{itemize}
  \end{columns}
\end{frame}

\section{正文结构}

\begin{frame}{一页只讲一个观点}
  \begin{block}{实用原则}
    每页先写结论，再放证据。不要把论文正文直接搬到幻灯片上。
  \end{block}

  \begin{exampleblock}{推荐结构}
    标题 = 观点；正文 = 3 个以内的支撑点；页脚进度条负责提示节奏。
  \end{exampleblock}
\end{frame}

\begin{frame}{一个简单表格}
  \begin{table}
    \centering
    \begin{tabular}{lll}
      \toprule
      场景 & 推荐做法 & 原因 \\
      \midrule
      答辩 & 少动画 & 降低不确定性 \\
      组会 & 多图表 & 方便讨论 \\
      技术分享 & 多代码 & 便于复现 \\
      \bottomrule
    \end{tabular}
  \end{table}
\end{frame}

\section{结束}

\begin{frame}{下一步}
  \begin{enumerate}
    \item 从主题列表里选一个风格
    \item 复制这份模板
    \item 先写大纲，再填每页内容
  \end{enumerate}
\end{frame}

\end{document}
